% \numberwithin{equation}{section}

\chapter{Resultats} \label{sec:resultats}
En aquesta secció es presenten els càlculs més importants i representatius del problema realitzats en els codis de l'Apèndix~\ref{ap.B}.
En l'anàlisi es segueix el mateix ordre que en les seccions anteriors: primerament es tracta el problema d'Alfvén i posteriorment el dels modes ràpids i lents.
Els valors numèrics emprats són els de la Taula~\ref{tab:valors}.
\begin{table}[H]
	\centering
	\begin{tabular}{|c|c|c|c|c|}
		\hline
		$c_{s c}^2$ & $c_{s p}^2$ & $v_{A c}^2$ & $v_{A p}^2$ & $k_z$ \\
		\hline
		$1$ & $6$ & $20$ & $250$ & $0.01$ \\
		\hline
	\end{tabular}
	\caption{Valors numèrics de velocitats adimensionals del sistema.
	El valor de la constant d'acoblament només s'empra en la Subsec.\ref{subsubsec:kz_constant}.}
	\label{tab:valors}
\end{table}
\vspace{-10pt}
Els valors numèrics són simulats a causa de certes limitacions computacionals:
per valors més realistes de les velocitats del so i d'Alfvén,% a la protuberància i la corona,
Dedalus dóna autofuncions i autovalors erronis.
Per contrastar els resultats amb valors realistes d'aquests paràmetres, es poden consultar els articles de \textcite{oscillation_modes_II} i~d'\textcite{oscillations_quiescent_solar_prominence}.

Un dels punts d'aquest treball és comparar els autovalors obtinguts per Dedalus amb aquells calculats analíticament a partir de combinacions sinusoïdals. % trigonomètriques
Per això es realitza en ambdós subproblemes el càlcul de l'error relatiu, $\eta$, de manera que
\begin{equation}
	\eta = \abs{\frac{\omega_D - \omega}{\omega}}, \tag{5.1} \label{eq.5.1}
\end{equation}
on $\omega$ és l'autovalor analític i $\omega_D$ l'obtingut per Dedalus.
Les expressions de les solucions matemàtiques es troben en major detall en l'article de \textcite{oscillation_modes_II}
i s'han calculat a partir d'una rutina numèrica.

Per determinar el comportament de cada mode es segueix una notació semblant a les Figs.~2~i~3 de l'article de \textcite{oscillation_modes_II}.
No obstant, en aquest treball s'ignora que els modes siguin ``kink'' (parells) o ``sausage'' (senars), per evitar una notació massa pesada.
A més, cada mode inverteix la seva paritat respecte el mode anterior i les autofuncions de $\tilde{v}_{x}$ i $v_z$ presenten sempre paritat oposada.

També s'afegeix la possibilitat que els modes siguin híbrids, no només interns o externs.
En l'article de \textcite{oscillation_modes_II} s'empra la paraula ``string'', tot i que més envant es redefineix el concepte de``mode híbrid'' en l'article d'\textcite{oscillations_quiescent_solar_prominence}.

En el cas dels modes d'Alfvén la única autofunció present és la que regeix el comportament del mode, $v_y$.
En canvi, en el cas considerat d'un acoblment feble, $k_z = 0.01$, en el problema dels modes ràpids i lents,
els modes lents tenen una amplitud dominant en la direcció paral·lela al camp, $\tilde{v}_{x}$,
mentre que en el ràpids governa l'amplitud en la direcció normal, $v_z$.

En resum, es consideren dues etiquetes pels modes d'Alfvén i tres pels modes ràpids i lents,
que determinen les diferents característiques de cada mode.
Es classifiquen d'acord amb les propietats de cada mode per $k_z$ = 0:
\begin{itemize}
	\item Segons l'estructura de l'autofunció dominant.
	\begin{itemize}
		\item \textbf{Fonamental (f)}: és el primer mode sempre, amb una única cresta i sense cap node.
		\item \textbf{Enèssims harmònics (nh)}: s'autofunció dominant es creua $n$ vegades amb l'eix d'abscisses i en total presenta $n + 1$ punts d'amplitud màxima i mínima, o crestes i valls, respectivament.
	\end{itemize}
	
	\item Segons les regions del sistema.
	\begin{itemize}
		\item \textbf{Externs (E)}: condicionats sobretot per la corona, la regió externa.
		\item \textbf{Interns (I)}: deguts principalment a les propietats de la protuberància, la zona interna.
		\item \textbf{Híbrids (H)}: influïts forta i simultàniament per la protuberància i la corona.
	\end{itemize}
	
	\item Segons la relació de dispersió.
	\begin{itemize}
		\item \textbf{Ràpids (R)}: presenten un creixement parabòl·lic de l'autovalor en funció de la constant d'acoblament.
		\item \textbf{Lents (L)}: amb una freqüència pràcticament constant a l'evolució del nombre d'ona,~$k_z$.
	\end{itemize}
\end{itemize}
Cal especificar que degut a la inexistència d'una relació de dispersió en el problema d'Alfvén, Eq.~(\ref{eq.3.30}),
no s'empra l'etiqueta per discriminar entre modes ràpids i lents.
Com ja s'ha comentat, aquesta classificació és desenvolupada en l'article d'\textcite{oscillations_quiescent_solar_prominence}.





\section{Modes d'Alfvén} \label{subsec:resultats_modes_Alfvén}
Com ja s'ha comentat, aquest cas s'assembla a una corda amb estructura longitudinal atès que representa un problema unidimensional,
tot i que la velocitat del mode és de l'estil~$v_y(x)$.
A més, no existeix cap tipus de relació de dispersió per aquesta situació degut a la independència dels autovalors respecte al nombre d'ona.

L'error relatiu dels autovalors resultants del problema d'Alfvén és considerablement petit.
Aquest creix a mesura que el nombre del mode augmenta, encara que no ho fa linealment.
Pels primers 50 autovalors el comportament de l'error és caòtic.
A partir d'aquí s'estabilitza i creix relativament poc, fins que a partir dels 200 creix exponencialment.
Aquest comportament es pot apreciar en la Fig.~\ref{fig:error_vA}.%figura que segueix.
\begin{figure}[H]
	\centering
	\includegraphics[width=0.915\textwidth]{Figures_Alfven_modes/eigenvalue_error_vA.png}
	\caption{Error relatiu, $\eta$, dels autovalors analítics i els calculats amb Dedalus en el cas del problema d'Alfvén.
	Els autovalors calculats amb Dedalus, $\omega_D$, estan calculats pel cas $k_z = 0.01$,
	tot i que s'ha comprovat que són pràcticament idèntics si es calculen amb $k_z = 0$.}
	\label{fig:error_vA}
\end{figure}
\vspace{-10pt}


Les autofuncions dels modes d'Alfvén són idèntiques a les d'una corda amb estructura longitudinal.
Per consegüent, l'estructura ve determinada pel seu índex de mode, veure Fig.~\ref{fig:modes_vA}.
\begin{figure}[H]
	\centering
	\includegraphics[width=0.915\textwidth]{Figures_Alfven_modes/modes_Alfven.png}
	\caption{Estructura de les primeres quatre autofuncions dels modes d'Alfvén amb el màxim normalitzat a 1.
	Les línies verticals grises determinen els límits de la protuberància, $[-x_p, x_p]$.}
	\label{fig:modes_vA}
\end{figure}
\vspace{-10pt}

El mode fonamental, f, té un únic màxim i cada harmònic presenta un extrem més que l'anterior.
Els dos primers harmònics, 1h i 2h, són externs i la seva estructura espacial es veu poc afectada per la presència de la protuberància.
En canvi, el tercer harmònic, 3h, és un mode intern i la seva distribució de $v_y(x)$ té màxims locals a les fronteres de la protuberància.
Aquests canvis de tendència precipitats són deguts precisament a l'entrada de l'ona dins la protuberància.

Les autofuncions estan fortament condicionades per les diferències en la velocitat d'Alfvén a la protuberància i a la corona, especificades en l'Eq.~(\ref{eq.3.34}).
A més, la consideració dels extrems fixes imposa solucions harmòniques, com s'ha comentat anteriorment.
Aquesta velocitat és molt major en l'interior, la protuberància, que en l'exterior, la corona,
de manera que la variació en velocitat dels modes en la protuberància és major, com s'aprecia en la Fig.~\ref{fig:modes_vA}.

Les parts imaginàries de les autofuncions són nul·les i no es presenten gràficament.
La informació rellevant és continguda en la part real.
En el càlcul numèric, aquestes parts imaginàries no són estrictament zero degut a la limitació de 16 xifres significatives de les variables float64 de Python,
tot i que són de l'ordre de $10^{-13}$ respecte a la part real normalitzada.

%Com ja s'ha esmentat, per advertir el comportament extern, intern o híbrid dels modes s'ha fet servir el quadern de Mathematica, que es pot trober en apèndix~[\ref{ap.A}].




\section{Modes ràpids i lents} \label{subsec:resultats_modes_r_i_l}
En aquest cas és convenient començar per la representació de la relació de dispersió.
És a dir, l'evolució dels autovalors, $\omega$, segons la constant d'acoblament, que no és més que el nombre d'ona en la direcció $\hat{z}$, $k_z$.
Llavors, d'acord amb les Eqs.~(\ref{eq.3.31})~i~(\ref{eq.3.32}) l'única direcció del vector d'ona que acobla els modes és la paral·lela a la protuberància, $k_z$.

\begin{figure}[H]
	\centering
	\includegraphics[width=0.95\textwidth]{Figures_F&S_modes/dd_modes_normals.png}
	\caption{Diagrama de dispersió.}
	\label{fig:dd}
\end{figure}
\vspace{-10pt}
De nou, les etiquetes corresponen a l'explicació anterior.
Tanmateix en aquest cas s'afegeix la tercera per indicar si el mode és ràpid, R, ò lent, L, d'acord amb l'article de \textcite{oscillation_modes_II}.
Aquesta qualitat s'aprecia directament amb l'evolució dels modes: els ràpids són parabòl·lics i els lents tenen una freqüència gairebé constant.
Cada color i etiqueta del diagrama de dispersió de la Fig.~\ref{fig:dd} correspon a un mode del problema.
Aquests són analitzats més envant en el treball i es representen les seves autofuncions amb els colors i etiquetes que els pertoquen.

Per comprovar la validesa dels autovalors calculats amb Dedalus, es calcula l'error relatiu, $\eta$, i es representa en la Fig.~\ref{fig:error_relatiu_r_i_l} per aquest cas.
La seva evolució és semblant a la del cas d'Alfvén de la Fig.~\ref{fig:error_vA}.
Ara bé, no és exactament igual. En aquest cas l'error és petit i anàrquic en els primers 50 autovalors,
s'estabilitza un poc més alluny, als 150, i als 200 es dispara.%, de nou per qüestions del càlcul computacional.
\begin{figure}[H]
	\centering
	\includegraphics[width=0.95\textwidth]{Figures_F&S_modes/eigenvalue_error.png}
	\caption{Error relatiu dels autovalors en el cas dels modes ràpids i lents.
	Els autovalors calculats amb Dedalus, $\omega_D$, estan calculats pel cas $k_z = 0.01$,
	tot i que s'ha comprovat que són pràcticament idèntics si es calculen amb $k_z = 0$.}
	\label{fig:error_relatiu_r_i_l}
\end{figure}
\vspace{-10pt}

Els modes ràpids i lents de la Fig.~\ref{fig:dd} resulten dels acoblaments dels diferents modes normals resultants de les Eqs.~(\ref{eq.3.31})~i~(\ref{eq.3.32}).
Per cada valor de kz, aquests s'indexen a partir de $n=1$. Els primers modes normals tenen una correspondència exacta amb els primers modes lents de la Fig.~\ref{fig:dd}.
Ara bé, el primer mode ràpid romp aquesta correspondència ja que presenta un canvi de comportament aprop de $k_z \sim 1$.
Aquesta situació d'intercanvi de comportament dels modes normals es defineix en l'article d'\textcite{oscillations_quiescent_solar_prominence}
com un ``encreuament evitat''.

Per entendre bé la situació es representa l'evolució del mode normal amb $n = 5$, inicialment de color morat, i s'analitzen els seus canvis a mida que s'incrementa $k_z$.
%S'ha escollit aquest mode atès que presenta dos canvis de tendència, així com el mode $n = 6$, de primeres representat en color grana. En qualsevol cas,
El cinquè mode normal pateix el primer encreuament evitat, aprop de $k_z \sim 1$.
Com es pot apreciar en les Figs.~\ref{fig:dd}~i~\ref{fig:dd_mode_5}, el mode $n = 5$ és inicialment lent (3hFL),
però aprop de $k_z \sim 1$ hi ha una inversió del comportament i passa a ser ràpid.
El creixement sobtat de l'autovalor ho determina, a diferència de la invariància d'aquest en el mode lent, en $k_z < 1$.
Envoltant de $k_z \sim 1.8$ es desfà el canvi i es recupera el mode lent (4hEL),
ara bé, és un mode lent harmònic d'un ordre superior a l'inicial.
\begin{figure}[H]
	\centering
	\includegraphics[width=0.85\textwidth]{Figures_F&S_modes/dd_mode_5.png}
	\caption{Evolució del mode normal amb $n = 5$.}
	\label{fig:dd_mode_5}
\end{figure}
\vspace{-10pt}

%Aquests fets es comproven més envant dins aquesta secció amb l'estudi de les autofuncions. 


L'evolució del mode normal $n = 5$ depèn del mode normal anterior, $n = 4$, de la Fig.~\ref{fig:dd}.
Aquest és inicialment un mode ràpid, amb un creixement parabòl·lic, fins a $k_z \sim 1$.
És en aquest punt on es realitza un encreuament evitat amb el mode superior degut a que aquests no poden solapar-se en el diagrama de dispersió.

Amb el fi d'observar el canvi de conducta dels modes,
en la Fig.~\ref{fig:dd_avoided_crossing} s'amplia el diagrama de dispersió de la Fig.~\ref{fig:dd} aprop del canvi, en $k_z \sim 1$.
\begin{figure}[H]
	\centering
	\includegraphics[width=0.85\textwidth]{Figures_F&S_modes/first_anticrossing.png}
	\caption{Encreuament evitat entre els modes fHR i 3hEL.}
	\label{fig:dd_avoided_crossing}
\end{figure}
\vspace{-10pt}
En aquest cas s'interpolen linealment els punts del diagrama de dispersió per demostrar explícitament que es tracta d'un encreuament evitat.
Així, és evident que els modes en sí no es creuen, sinó que hi ha un intercanvi de tendència d'aquests, així com s'havia explicat anteriorment.

%---------------------------------------------------------------------------------------------------------------------------------------------------
%Vz va en direcció del camp i no li afecta molt l'acoblament de kz --- afegir al tfg --- a quina secció? teoria? i despres s'observa en l'evolució del mode fHR --- sí!

%La velocitat paral·lela a la direcció d'acoblament, $v_z$, no es veu molt alterada per l'acoblament precisament.
%Aquest fet és un reflexe de les equacions dels modes ràpids i lents, Eqs.~(\ref{eq.3.31}) i (\ref{eq.3.32})
%---------------------------------------------------------------------------------------------------------------------------------------------------




%\subsubsection{Autofuncions per $k_z$ fixa.}
\subsection{Autofuncions amb constant d'acoblament fixa, $k_z = 0.01$} \label{subsubsec:kz_constant}

Una vegada estudiat el diagrama de dispersió es representen les autofuncions, és a dir, l'estructura dels modes ràpids i lents al llarg de l'eix $\hat{x}$.
A la Fig.~\ref{fig:autofuncions_reals_individuals} es considera una constant d'acoblament petita, $k_z = 0.01$,
per poder apreciar la diferència amb els primers quatre modes d'Alfvén de la Fig.~\ref{fig:modes_vA} de l'apartat anterior.
\begin{figure}[H]
	\centering
	\includegraphics[width=0.95\textwidth]{Figures_F&S_modes/Autofuncions_reals_individuals.png}
	\caption{Estructura de les primeres quatre autofuncions amb els màxims normalitzats a 1, associades als modes ràpids i lents per $k_z = 0.01$.
	Les línies verticals grises determinen els límits de la protuberància, $[-x_p, x_p]$.}
	\label{fig:autofuncions_reals_individuals}
\end{figure}
\vspace{-10pt}
Cada mode es correspon a l'etiqueta i color pertinent de la Fig.~\ref{fig:dd}. 
Els d'ordre superior no es respresenten per qüestions visuals, tot i que s'han treballat, com demostren les etiquetes.

El canvi més remarcable respecte als modes d'Alfvén de la Fig.~\ref{fig:modes_vA} és que cada mode, ja sigui ràpid o lent,
presenta dues autofuncions: $\tilde{v}_{x}$ i $v_z$.
En els tres primers modes predomina el terme de $\tilde{v}_{x}$, llavors tenen la mateixa estructura que els primers tres modes d'Alfvén.
En canvi, en el quart mode es giren les tornes i l'autofunció predominant és ${v}_{z}$.
Això indica que les tres primeres solucions corresponen a modes lents mentre que la quarta a un ràpid, en el cas de $k_z = 0.01$.
Per tant, la situació està en acord amb el diagrama de dispersió de la Fig.~\ref{fig:dd},
segons el qual, amb un valor de la constant d'acoblament major, el quart mode també ha de ser lent.

Una altra lectura que ofereix la representació de les autofuncions %dels modes ràpids i lents
de la Fig.~\ref{fig:autofuncions_reals_individuals}
és que amb una constant, $k_z$, molt petita l'acoblament és certament feble i l'autofunció no dominant és pràcticament nul·la.
En el límit $k_z \rightarrow 0$, aquesta autofunció no dominant és igual a zero $\forall x$.

En la següent subsecció es tracta amb detall l'evolució dels modes segons segons la importància de l'acoblament.
Per aquest motiu no es representen les autofuncions d'altres valors de la constant $k_z$.



%---------------------------------------------------------------------------------------------------------------------------------------------------
% DEMANAR A MONXO SI LI PUC ENXUFAR AQUÍ UNA FIGURA MÉS PER COMPROVAR AIXÒ QUE DIC!
%---------------------------------------------------------------------------------------------------------------------------------------------------






\subsection{Evolució de les autofuncions en funció de l'acoblament, $k_z$} \label{subsubsec:autofuncions_kz}
Per apreciar el canvi de comportament es comparen les autofuncions aprop de l'encreuament evitat representat en la Fig.~\ref{fig:dd_avoided_crossing}.
Llavors es representa l'estructura dels modes entorn al valor de la constant d'acoblament de $k_z \sim 1$ en la Fig.~\ref{fig:autofuncions_anticreuament}.
\begin{figure}[H]
	\centering
	\includegraphics[width=0.925\textwidth]{Figures_F&S_modes/Evolució_autofuncions_n_4_subplots.png}
	\caption{Evolució de les autofuncions normalitzades a 1 del mode fHR en el primer encreuament evitat.
	Les línies verticals grises marquen els límits de la protuberància, $[-x_p, x_p]$.}
	\label{fig:autofuncions_anticreuament}
\end{figure}
\vspace{-10pt}
En aquest cas es veu el canvi de dominància de l'autofunció perpendicular al camp, $v_z$, enfront a l'autofunció paral·lela a aquest, $\tilde{v}_{x}$.

Per $k_z = 0.90$ el mode representat en vermell encara és completament dominat per la velocitat en direcció paral·lela al vector d'ona, $v_z$.
En canvi, així com augmenta la constant d'acoblament aquesta, $v_z$, minva i la velocitat paral·lela al camp, $\tilde{v}_{x}$, creix,
segons amb la normalització explicada amb anterioritat.

Finalment, en $k_z = 1.10$ el mode representat en rosa evidencia que la velocitat $\tilde{v}_{x}$ és major que $v_z$.
Per tant, s'ha duit a terme un intercanvi de comportament dominant en el mode normal, de ràpid a lent,
així com s'havia predit just abans de començar la~Subsec.~\ref{subsubsec:autofuncions_kz}.
%\hyperref[subsubsec:autofuncions_kz]{Subsec.~5.2.1}.
%així com s'havia predit en el final de la \hyperref[subsec:resultats_modes_r_i_l]{Subsec.~5.2}.

Per acabar s'estudia l'evolució dels dos modes fonamentals híbrids, fHL i fHR,
en funció del valor de l'acoblament, $k_z$, en un rang ampli de valors d'aquest paràmetre, sempre situats enfora dels encreuaments evitats pel mode fHR.

Inicialment, a la Fig~\ref{fig:fHL} s'analitza el mode híbrid lent, fHL, en vermell en el diagrama de dispersió de la~Fig.~\ref{fig:dd}.
\begin{figure}[H]
	\centering
	\includegraphics[width=0.95\textwidth]{Figures_F&S_modes/Evolució_autofuncions_fHL_subplots.png}
	\caption{Evolució de les autofuncions amb màxims normalitzats a 1 del mode fHL segons els encreuaments evitats.
	Les línies verticals grises marquen els límits de la protuberància, $[-x_p, x_p]$.}
	\label{fig:fHL}
\end{figure}
\vspace{-10pt}
Les autofuncions $\tilde{v}_{x}$ són pràcticament invariants, independents de la constant d'acoblament.
Aquest fet és d'esperar ja que en la relació de dispersió el mode és lent, amb freqüència quasi constant, i no es troba amb cap ràpid.
En conseqüència, les autofuncions es solapen atès que són pràcticament idèntiques.
Amb una gran ampliació de la Fig.~\ref{fig:fHL}a en el màxim es poden apreciar les diferències, tot i que es perd la forma de l'autofunció $\tilde{v}_{x}$.

Es requereix l'ús d'un rang de l'eix vertical diferent per l'autofunció no dominant, $v_z$,
degut a que si es representa a la mateixa escala que la dominant, $\tilde{v}_{x}$, es perd la informació i sembla una línia recta,
com en el cas de la Fig.~\ref{fig:autofuncions_reals_individuals}.
En aquell cas s'ha prioritzat un rang de l'eix vertical gran per poder apreciar l'estructura de l'autofunció dominant.

La màxima amplitud de l'autofunció no dominant, $v_z$,
és de l'ordre de $\sim 0.005$ respecte al major màxim d'amplitud de l'autofunció dominant, $\tilde{v}_{x}$.

Això mateix passa amb l'autofunció vermella en la pròpia Fig.~\ref{fig:fHL}b,
el mode no dominant és tan inferior al dominant en valors de $k_z$ baixos, que el cas de $k_z = 0.01$ l'autofunció sembla constant,
tot i que té la mateixa estructura que la resta de casos representats, del tipus $-\sin (\pi x)$.
D'aquesta manera es reflecteix la gran importància de la constant d'acoblament.
Quan és nul·la els modes estan desacoblats i es recupera el cas d'Alfvén, com prediuen les Eqs.~(\ref{eq.3.31}) i (\ref{eq.3.32}).



El segon cas estudiat és el mode híbrid ràpid, fHR, en groc en la relació de dispersió de la Fig.~\ref{fig:dd}.
Tenint en compte que es tracta d'un mode ràpid, el seu autovalor, $\omega$, creix de manera parabòl·lica respecte a la constant d'acoblament, $k_z$.
Llavors s'espera que les autofuncions es veguin modificades i reflecteixin l'evolució de l'autovalor, veure la Fig.~\ref{fig:fHR}
\begin{figure}[H]
	\centering
	\includegraphics[width=0.95\textwidth]{Figures_F&S_modes/Evolució_autofuncions_fHR_subplots.png}
	\caption{Evolució de les autofuncions amb màxims normalitzats a 1 del mode fHR segons els encreuaments evitats.
	Les línies verticals grises marquen els límits de la protuberància, $[-x_p, x_p]$.}
	\label{fig:fHR}
\end{figure}
\vspace{-10pt}
Efectivament en aquest cas les autofuncions sí pateixen un canvi estructural degut al canvi de $k_z$.
El que més destaca és l'evolució de les autofuncions no dominants, $\tilde{v}_{x}$.
Per començar, les amplituds màximes dels modes creixen a major acoblament, com en el cas anterior.
Cada vegada que es duu a terme un encreuament evitat, els modes de les autofuncions no dominants retenen la longitud d'ona del mode lent amb el que han evitat creuar-se.
És per aquest motiu que l'estructura, $\tilde{v}_{x}$, varia considerablement a cada autofunció representada, ja que aquesta s'ha calculat després d'un nou encreuament evitat.
D'aquesta manera, a mesura que l'acoblament augmenta, $\tilde{v}_{x}$ conté longituds d'ona cada vegada més petites.

En aquest cas es pot apreciar també com amb poc acoblament, $k_z = 0.01$,
corba vermella, la dominància de l'autofunció $v_z$ és encara major precisament perque els modes es troben poc lligats.
A més, s'intueix que es tracta del mateix mode que els altres però amb un acomblament major,
ja que $\tilde{v}_{x}$ creua l'eix de les abscisses tres vegades, com la resta de corbes.
A causa de l'escala que pràcticament sembla una recta, tot i que amb una fina mirada s'aprecia que no ho és en cap cas.

El major màxim de les autofuncions no dominants és $\sim 0.075$,
unes 15 vegades major que el major màxim de les autofuncions no dominants del primer mode híbrid lent, fHL, de la Fig.~\ref{fig:fHL}b.
A més, els dos màxims de les no dominants tenen lloc en el cas de major acoblament, amb $k_z = 2.5$.
Llavors es dedueix que la lligadura té un efecte major en els modes ràpids que en els lents.
En els cas de $k_z = 0.75$ la diferència és encara major ja que el màxim del no dominant del mode fHL, $v_z$, és de $\sim 0.0025$,
mentre que el del no dominant del mode fHR, $\tilde{v}_{x}$, és de $\sim 0.06$, unes 24 vegades major.
Definitivament l'acoblament té una major influència en l'estructura dels modes ràpids que en els lents, tot i que no és constant.

Com que es tracta d'un mode ràpid, les autofuncions dominants són les de $v_z$. Aquestes conserven l'estructura bàsica de mode fonamental, f,
i presenten lleugeres modificacions així com augmenta $k_z$ a causa d'un major acoblament amb les autofuncions no dominants, $\tilde{v}_{x}$.
El seu comportament es veu alterat com a un estrenyiment de la campana que conforma l'autofunció dominant a mesura que l'acoblament creix.

L'etiqueta de fonamental, f, en els modes fHL i fHR es deu a l'estructura de l'autofunció dominant del mode corresponent.
Es comprova en els dos casos que les autofuncions presenten una paritat oposada, llavors les autofuncions no dominants són sempre harmònics senars.
%, 1h i 3h respectivament.
